% =============================================================================
% DRAFT: Revised Section 3.3 and Section 7 for Paper 4A (preprint)
%
% Replaces:
%   - Current Sec 3.3 ("DLN stage index and three discrete observer-quality states")
%   - Current Sec 7  ("DLN-series connection and stage mapping")
%
% Design principle: DLN formal objects (G, M, R) are IMPORTED from Paper 1
% and SPECIALIZED to QD/SBS.  No redefinition.
% =============================================================================


% =====================================================================
%  NEW SECTION 3.3
% =====================================================================

\subsection{Observer inference topology: DLN instantiation in quantum Darwinism}
\label{sec:qstates}

The resource triple $Q_O=(R_O,\Lambda_O,\tau_O)$ parameterizes observer
\emph{limitations}---how many fragments can be accessed, how each is
distorted, and how long coherence can be maintained.
It does not parameterize the observer's \emph{inference structure}: the
topology by which evidence from multiple fragments is organized, combined,
and---when necessary---reorganized.

The Dot--Linear--Network (DLN) framework
\citep{WuDLNCompression2026preprint} provides a graph-theoretic
formalization of exactly this missing variable.
DLN characterizes cognitive or information-processing stages by two
objects:
a \emph{belief-dependency graph}~$G$ governing within-episode inference,
and a \emph{revision graph}~$\mathcal{R}$ over a model
space~$\mathcal{M}$ governing when and how the active inference
structure can be changed.
We now show that both objects instantiate directly in the QD/SBS setting:
$G$ as the conditional-independence structure of the SBS state, and
$\mathcal{R}$ as the set of physically realizable transitions between
measurement strategies.


\subsubsection{SBS states as bipartite factor DAGs (Level~1: belief graph $G$)}

\begin{proposition}[SBS conditional independence as a bipartite DAG]
\label{prop:sbs-dag}
Let $\sigma_{S\cE}$ be an SBS state (Definition~1) with pointer variable
$X$ taking values in~$\cX$.  The tensor-product form
\begin{equation}
  \sigma_{S\cE}
  = \sum_{x\in\cX} p_x \ket{x}\!\bra{x}_S
    \otimes \bigotimes_{k=1}^{N} \sigma^{(x)}_{\cE_k}
  \label{eq:sbs-factor}
\end{equation}
defines a bipartite directed acyclic graph (DAG): a single latent node
$X$ with directed edges $X \to \cE_k$ for $k=1,\dots,N$, and no edges
among fragment nodes.
Conditional on $X=x$, the fragment states
$\{\sigma^{(x)}_{\cE_k}\}_{k=1}^{N}$ are mutually independent:
\[
  \cE_j \;\perp\!\!\!\perp\; \cE_k \mid X
  \qquad \text{for all } j\neq k.
\]
\end{proposition}

\begin{proof}
The product structure
$\bigotimes_{k=1}^{N}\sigma^{(x)}_{\cE_k}$ conditional on $X=x$ is the
defining property of conditional independence in a tensor-product state.
In the classical reduction (diagonal in the pointer basis, with each
fragment traced to measurement outcomes $y_k$), the joint distribution
factorizes as
$p(x,y_1,\dots,y_N)=p(x)\prod_k p(y_k\mid x)$,
which satisfies the Markov property of the stated DAG:
the pointer node $X$ d-separates all fragment nodes
\citep{Pearl1988}.
\end{proof}

\begin{remark}[SBS as a DLN factor graph]
\label{rem:sbs-factor}
The DAG of Proposition~\ref{prop:sbs-dag} is precisely the
\emph{belief-dependency graph}~$G$ of a Network-stage agent in the DLN
framework \citep[Sec.~2.1]{WuDLNCompression2026preprint}: a bipartite
graph with $F$ shared latent factors connected to $K$ evidence nodes.
In the QD setting, $F=|\cX|$ (pointer alphabet size, typically
$|\cX|=2$) and $K=N$ (number of environment fragments).
Since QD objectivity requires high redundancy, i.e.\ $N\gg 1$, the
condition $F \ll K$ that drives the DLN compression thesis is
automatically and generically satisfied in any QD scenario.
\end{remark}


\subsubsection{Decoder classes as DLN stages}

We define three nested classes of decoders on the $m$-fragment register,
directly instantiating the DLN stage classification of
\citep[Sec.~2.1]{WuDLNCompression2026preprint} in the QD measurement
setting.

\begin{definition}[Decoder classes (DLN stages in QD)]
\label{def:q-states}
Let $A\subseteq\{1,\dots,N\}$ with $|A|=m$, and let
$\rho_A^{(x)} := \bigotimes_{k\in A}\Lambda_O(\sigma^{(x)}_{\cE_k})$
denote the calibrated conditional states (SBS witness experiment,
$\cX=\{0,1\}$).
A \emph{decoder} is a measurement--decision procedure producing an
estimate $\widehat{X}\in\cX$ from the $m$-fragment register.
We define three decoder classes, each corresponding to a DLN stage
\citep[Sec.~2.1]{WuDLNCompression2026preprint}:
%
\begin{enumerate}[leftmargin=2.2em,label=(\roman*)]
\item \textbf{$q_{\mathrm{D}}$ (Dot): memoryless decoding.}
The decoder acts on at most one fragment $\cE_k$ at a time and retains
no state across fragments.
The belief-dependency graph is empty ($G=\emptyset$): no evidence
propagates between fragment nodes.

\item \textbf{$q_{\mathrm{L}}$ (Linear): product decoding with
classical postprocessing.}
The decoder applies a product POVM $M=\bigotimes_{k\in A}M_k$, followed
by an arbitrary classical decision rule on the $m$~outcomes.
The belief-dependency graph is a null graph on $m$ evidence nodes: each
fragment is processed independently, and learning from fragment~$j$ does
not update inference about fragment~$k$.

\item \textbf{$q_{\mathrm{N}}$ (Network): collective decoding.}
The decoder may apply an arbitrary joint POVM on the full register
$\bigotimes_{k\in A}\cE_k$.
The belief-dependency graph is the bipartite DAG of
Proposition~\ref{prop:sbs-dag}: the observer's inference explicitly
represents the pointer~$X$ as a shared latent parent of all fragment
nodes, and evidence from any single fragment propagates through~$X$ to
update predictions for all others.
\end{enumerate}
\end{definition}


\begin{proposition}[Strict nesting of decoder classes]
\label{prop:nesting}
Let $\mathsf{Dec}_D$, $\mathsf{Dec}_L$, $\mathsf{Dec}_N$ denote the
sets of all decoders in classes $q_{\mathrm{D}}$, $q_{\mathrm{L}}$,
$q_{\mathrm{N}}$, respectively.  Then
\[
  \mathsf{Dec}_D \;\subsetneq\; \mathsf{Dec}_L
                  \;\subsetneq\; \mathsf{Dec}_N.
\]
Both inclusions are strict whenever $m\ge 2$ and the conditional
fragment states $\{\sigma^{(x)}_{\cE_k}\}$ are not identical
across~$x$.
\end{proposition}

\begin{proof}
A single-fragment measurement on $\cE_k$ is the product POVM with
$M_j=\mathbf{1}$ for $j\neq k$, establishing
$\mathsf{Dec}_D\subseteq\mathsf{Dec}_L$.
Any product POVM is a special case of a joint POVM, giving
$\mathsf{Dec}_L\subseteq\mathsf{Dec}_N$.
For strictness of the first inclusion: any POVM acting on two or more
fragments is in $\mathsf{Dec}_L\setminus\mathsf{Dec}_D$.
For strictness of the second: the Helstrom-optimal measurement on
$\rho_A^{(0)}$ versus $\rho_A^{(1)}$ is generically entangled across
fragments when the per-fragment overlaps $c_k$ are not~$0$ or~$1$
\citep{Helstrom1976}, placing it in
$\mathsf{Dec}_N\setminus\mathsf{Dec}_L$.
\end{proof}


\subsubsection{Quantitative content: stage-dependent error exponents and redundancy}

The strict nesting of decoder classes has quantitative consequences for
the central QD observable: the number of fragments required to achieve a
given decoding error.

\begin{proposition}[Tight exponent separation across DLN stages, pure-state records]
\label{prop:tight-gap}
Let $\ket{\psi_0},\ket{\psi_1}$ be pure states with overlap
$c:=|\braket{\psi_0}{\psi_1}|\in(0,1)$.
Consider the task of discriminating $\ket{\psi_0}^{\otimes m}$ from
$\ket{\psi_1}^{\otimes m}$ under equal priors.
The optimal error exponents achievable by each decoder class are:
%
\begin{align}
  \xi_{\mathrm{N}}
    &:= \lim_{m\to\infty}
        -\tfrac{1}{m}\log P_e^{\mathrm{coll}}(m)
    = -\log(c^2),
  \label{eq:xi-N}\\[4pt]
  \xi_{\mathrm{L}}
    &:= \lim_{m\to\infty}
        -\tfrac{1}{m}\log P_e^{\mathrm{prod}}(m)
    = -\log(c),
  \label{eq:xi-L}\\[4pt]
  \xi_{\mathrm{D}}
    &:= 0\qquad
    \bigl(\text{$P_e^{\mathrm{single}}
    = \tfrac{1}{2}(1-\sqrt{1-c^2})$,
    independent of $m$}\bigr).
  \label{eq:xi-D}
\end{align}
%
In particular:
\begin{enumerate}[leftmargin=2em,label=(\alph*)]
\item The exponent ratio between Network and Linear stages is
      exactly~$2$:
      $\xi_{\mathrm{N}}/\xi_{\mathrm{L}} = 2$.
\item Both exponents are \emph{achievable}:
      $\xi_{\mathrm{N}}$ by the Helstrom measurement on the joint
      $m$-copy register \citep{Audenaert2007,NussbaumSzkola2009},
      and $\xi_{\mathrm{L}}$ by the optimal single-copy POVM
      (which induces classical distributions with Bhattacharyya
      coefficient $B=c$ for pure states \citep{FuchsCaves1994})
      followed by the Neyman--Pearson likelihood-ratio test on the
      i.i.d.\ classical outcomes.
\item The Dot stage has zero exponent: a memoryless decoder's error
      is bounded below by the single-copy Helstrom error regardless
      of $m$, because no evidence is accumulated.
\end{enumerate}
\end{proposition}

\begin{proof}
Equation~\eqref{eq:xi-N} is the quantum Chernoff exponent for pure
states \citep{Audenaert2007,NussbaumSzkola2009}.
For \eqref{eq:xi-L}: any single-copy POVM on pure states
$\ket{\psi_0},\ket{\psi_1}$ induces classical distributions $p,q$ with
Bhattacharyya coefficient
$B(p,q)=\sum_y\sqrt{p(y)q(y)}\ge c$
\citep{FuchsCaves1994,Jozsa1994}.
For the classical Chernoff coefficient,
$\min_{s\in[0,1]}\sum_y p(y)^s q(y)^{1-s} \ge B$ (since $s=\tfrac12$
is feasible and yields~$B$), so the $m$-copy product-measurement
coefficient is at least $c^m$, giving exponent at most $-\log c$.
Achievability: the Helstrom single-copy POVM on pure states with
overlap~$c$ induces classical distributions with $B = c$ exactly
(direct calculation of the binary Bhattacharyya coefficient for the
Helstrom outcome probabilities
$p_\pm = \tfrac{1}{2}(1\pm\sqrt{1-c^2})$), so the classical Chernoff
exponent of the induced i.i.d.\ channel is precisely $-\log c$.
Equation~\eqref{eq:xi-D}: a Dot decoder accesses one fragment per
episode with no memory, so its error equals the single-copy Helstrom
error irrespective of the total number of available fragments.
\end{proof}


\begin{remark}[Mixed-state fragment records]
\label{rem:mixed-state-gap}
Proposition~\ref{prop:tight-gap} is stated for pure-state records,
which arise naturally in the central-spin worked example
(Sec.~\ref{sec:centralspin}) and in any pure-dephasing model.
For mixed-state conditional fragment states, the exponent ratio
$\xi_{\mathrm{N}}/\xi_{\mathrm{L}}$ lies in $[1,2]$:
the lower bound~$1$ is saturated when the conditional states commute
(classical records, where collective measurements offer no advantage
over product measurements),
and the upper bound~$2$ is saturated in the pure-state limit.
The ratio is a monotone function of the non-commutativity of the
conditional fragment states, providing a continuous interpolation between
classical ($\xi_{\mathrm{N}}/\xi_{\mathrm{L}}=1$) and fully quantum
($\xi_{\mathrm{N}}/\xi_{\mathrm{L}}=2$) records.
A complete characterization of the mixed-state exponent gap in terms of
the fragment state geometry (e.g., as a function of quantum
R\'{e}nyi divergences \citep{MosonyiOgawa2015}) is an open problem.
\end{remark}


\begin{corollary}[Stage-dependent redundancy requirement]
\label{cor:redundancy-stages}
Fix a target error $\delta\in(0,\tfrac12)$ and suppose all $N$ fragments
carry identical pure-state records with overlap~$c$.
Let $m_q(\delta)$ denote the minimum number of fragments sufficient for
a stage-$q$ observer to achieve error $\le\delta$.
To leading order in $\log(1/\delta)$:
\[
  m_{\mathrm{D}}=\infty\;\text{(if $P_e^{\mathrm{single}}>\delta$)},
  \qquad
  m_{\mathrm{L}}(\delta) \sim \frac{\log(1/(2\delta))}{-\log c},
  \qquad
  m_{\mathrm{N}}(\delta) \sim \frac{\log(1/(2\delta))}{-\log(c^2)}
  = \frac{m_{\mathrm{L}}(\delta)}{2}.
\]
A Network-stage observer requires half as many fragments as a
Linear-stage observer for the same error target.
\end{corollary}

\begin{proof}
Immediate from Proposition~\ref{prop:tight-gap} and the inversion
$m_q(\delta) \sim \log(1/(2\delta)) / \xi_q$.
\end{proof}


\begin{remark}[Redundancy ratio as a DLN compression constant]
\label{rem:compression}
In the language of \citep{WuDLNCompression2026preprint}, the ratio
$m_{\mathrm{L}}/m_{\mathrm{N}}=2$ is the \emph{compression efficiency
gain}: the reduction in evidence units required when the observer
exploits shared latent structure (the pointer variable) versus treating
evidence units independently.
This is the QD specialization of the $O(F)$-versus-$O(K)$ cost-scaling
thesis of \citep[Proposition~1(i)]{WuDLNCompression2026preprint}.
The ratio is an exact constant (not merely an asymptotic scaling
relation) because the pointer alphabet size $|\cX|$ is fixed while the
number of fragments~$N$ grows, placing QD in the regime $F\ll K$ where
the DLN compression advantage is maximal.
\end{remark}


\subsubsection{Model space and revision graph
              (Level~2: measurement-strategy transitions)}
\label{sec:revision-qd}

The DLN framework formalizes not only the within-episode belief structure
(the graph~$G$) but also the observer's capacity to \emph{revise} that
structure.
Following \citep[Sec.~2.3]{WuDLNCompression2026preprint}, this is
captured by a model space~$\mathcal{M}$ (the set of candidate belief
structures) and a revision graph $\mathcal{R}=(\mathcal{M},\mathcal{T})$
(a directed graph whose edges are the available transitions between
structures).
We now specialize these objects to the QD setting.

\begin{definition}[Measurement model space]
\label{def:model-space-qd}
For an observer $O$ with resource triple
$Q_O=(R_O,\Lambda_O,\tau_O)$, define the
\emph{measurement model space}
\[
  \mathcal{M}_O \;\subseteq\;
  \{\mathsf{Dec}_D,\;\mathsf{Dec}_L,\;\mathsf{Dec}_N\}
\]
as the set of decoder classes that are physically realizable
given~$Q_O$.
A decoder class $\mathsf{Dec}_q$ is in~$\mathcal{M}_O$ if and only if
the observer possesses the physical resources to implement at least one
decoder in that class:
\begin{itemize}
\item $\mathsf{Dec}_D\in\mathcal{M}_O$ always
  (memoryless decoding requires no inter-fragment coherence).
\item $\mathsf{Dec}_L\in\mathcal{M}_O$ whenever $\tau_O$ is sufficient
  to sequentially acquire $m$ fragment outcomes and store classical
  records.
\item $\mathsf{Dec}_N\in\mathcal{M}_O$ whenever $\tau_O$ is sufficient
  to maintain quantum coherence across the $m$-fragment register for
  the full acquisition-plus-decoding interval.
\end{itemize}
\end{definition}


\begin{definition}[Revision graph in QD]
\label{def:revision-qd}
The \emph{revision graph}
$\mathcal{R}_O = (\mathcal{M}_O,\,\mathcal{T}_O)$
is a directed graph over the observer's model space, where
$\mathcal{T}_O \subseteq \mathcal{M}_O\times\mathcal{M}_O$ is the set
of transitions the observer can execute.
Following the DLN classification
\citep[Sec.~2.3]{WuDLNCompression2026preprint}:
%
\begin{itemize}
\item \textbf{Dot.}
  $\mathcal{M}_O=\{\mathsf{Dec}_D\}$;
  $\mathcal{R}_O$ is a single vertex with no edges.
  No revision capacity.

\item \textbf{Linear.}
  $\mathcal{M}_O=\{\mathsf{Dec}_L\}$;
  $\mathcal{R}_O$ is a single vertex.
  The observer occupies a fixed point in model space.

\item \textbf{Network (expand-only).}
  $\mathcal{M}_O=\{\mathsf{Dec}_N,\mathsf{Dec}_L\}$ with a single
  directed transition
  $\mathsf{Dec}_N\to\mathsf{Dec}_L$
  (degradation when coherence drops below the collective-decoding
  threshold).
  No return transition.

\item \textbf{Network (full cycle).}
  Both transitions
  $\mathsf{Dec}_N \rightleftharpoons \mathsf{Dec}_L$
  are available.
  $\mathcal{R}_O$ contains a directed cycle: the observer can degrade to
  product measurements when coherence is lost and recover collective
  measurements when coherence is restored.
\end{itemize}
\end{definition}


\begin{remark}[Physical mechanisms governing $\mathcal{R}$-transitions]
\label{rem:transition-physics}
Transitions in $\mathcal{R}_O$ are governed by the relation between the
observer's temporal horizon $\tau_O$ and the coherence time required for
collective decoding.
In laboratory settings, the relevant physical mechanisms include:
(i)~fluctuating dephasing rates (e.g., temperature-dependent $T_2$ in
spin-bath experiments \citep{Schlosshauer2007}),
(ii)~time-varying calibration noise (drift in detector efficiency or
alignment),
(iii)~changes in pointer stability (system relaxation modifying the SBS
quality parameter~$\varepsilon$).
An observer that monitors its own decoding performance---a Level~2
operation in the sense of
\citep[Sec.~2.3]{WuDLNCompression2026preprint}---and switches decoder
class in response, is executing a transition in~$\mathcal{R}_O$.
\end{remark}


\begin{proposition}[Bounded recovery under measurement-strategy revision]
\label{prop:recovery-qd}
Let an observer~$O$ initially operate in $\mathsf{Dec}_N$ (collective
decoding, exponent $\xi_{\mathrm{N}}=-\log(c^2)$) and suppose
coherence conditions degrade at time~$t_0$, forcing a transition to
$\mathsf{Dec}_L$ (product decoding, exponent $\xi_{\mathrm{L}}=-\log c$).
Suppose coherence conditions subsequently recover at time $t_1>t_0$.
\begin{enumerate}[label=(\roman*)]
\item If $\mathcal{R}_O$ contains the return transition
  $\mathsf{Dec}_L\to\mathsf{Dec}_N$ (full-cycle Network observer), the
  observer can recover the exponent $\xi_{\mathrm{N}}$ within a bounded
  verification interval after~$t_1$.
  The recovery time is determined by the monitoring protocol (Level~2
  overhead), analogous to the contraction bound in
  \citep[Proposition~1(iii)]{WuDLNCompression2026preprint}.
\item If $\mathcal{R}_O$ lacks the return transition (expand-only Network
  observer), the observer remains at exponent $\xi_{\mathrm{L}}$
  permanently, even after coherence is restored.
\end{enumerate}
The performance cost of being locked in $\mathsf{Dec}_L$ is a factor of~$2$
in the required fragment count (Corollary~\ref{cor:redundancy-stages}):
\[
  m_{\mathrm{L}}(\delta) = 2\,m_{\mathrm{N}}(\delta)
  \qquad\text{(pure-state records).}
\]
\end{proposition}

\begin{proof}
Part~(i): after coherence recovery, the conditions for
$\mathsf{Dec}_N\in\mathcal{M}_O$ are again satisfied.
A full-cycle observer with a Level~2 monitoring protocol
(cf.\ the shadow-model verification of
\citep[Sec.~5.1]{WuDLNCompression2026preprint}) detects that collective
decoding is again feasible and executes
$\mathsf{Dec}_L\to\mathsf{Dec}_N$.
The verification interval is bounded by the monitoring window length plus
a confirmation period, as established for the analogous contraction
mechanism in
\citep[Proposition~1(iii)]{WuDLNCompression2026preprint}.
Part~(ii): without the return transition, the observer has no mechanism
to re-enter $\mathsf{Dec}_N$, regardless of restored physical capacity.
The fragment-count penalty follows from
Corollary~\ref{cor:redundancy-stages}.
\end{proof}


\paragraph{Concrete monitoring protocol: coherence-gated decoder switching.}
\label{par:monitoring-protocol}
We now specify a physically realizable Level~2 protocol for executing
$\mathcal{R}$-transitions.
The observer partitions its accessible fragments into a \emph{decoding
set} $A_{\mathrm{dec}}$ and a \emph{monitoring set}
$A_{\mathrm{mon}}$, with $|A_{\mathrm{mon}}|=\lceil f\cdot m\rceil$
for a fixed monitoring fraction $f\in(0,1)$ (typically $f\approx 0.1$--$0.2$).
The monitoring set is reserved for Level~2 diagnostics and does not
contribute to the primary pointer estimate.

\begin{enumerate}[leftmargin=2em,label=\textbf{L2.\arabic*}]
\item \textbf{Cross-validation scoring (continuous).}
  In each observation episode:
  (a)~apply the current decoder (collective or product) to
  $A_{\mathrm{dec}}$, obtaining a pointer estimate $\widehat{x}$;
  (b)~apply independent single-fragment measurements to each
  $\cE_k\in A_{\mathrm{mon}}$, obtaining per-fragment estimates
  $\{\widehat{x}_k\}$;
  (c)~compute the \emph{consistency score}
  $s_t := |A_{\mathrm{mon}}|^{-1}\sum_{k\in A_{\mathrm{mon}}}
  \mathbf{1}[\widehat{x}_k = \widehat{x}]$.
  Maintain a rolling mean $\bar{s}_W$ over a window of the last $W$
  episodes.

\item \textbf{Expansion trigger
  ($\mathsf{Dec}_N\to\mathsf{Dec}_L$).}
  If $\bar{s}_W$ falls below the \emph{product-decoder baseline}
  $s_{\mathrm{prod}} := 1-P_e^{\mathrm{single}}$ (the consistency
  expected when $\widehat{x}$ is no more reliable than majority vote
  over monitoring fragments), the collective decoder is no longer adding
  value beyond what product decoding achieves.
  The observer executes
  $\mathsf{Dec}_N\to\mathsf{Dec}_L$.

  \emph{Physical trigger:} equivalently, the observer may directly
  monitor its coherence time via a Ramsey-contrast or spin-echo sequence
  on an ancilla qubit coupled to the fragment register.
  If the estimated $\widehat{T}_2$ drops below the collective-decoding
  threshold $\tau_{\mathrm{coll}}:= m\,t_{\mathrm{meas}}+t_{\mathrm{decode}}$,
  the transition is triggered.

\item \textbf{Contraction trigger
  ($\mathsf{Dec}_L\to\mathsf{Dec}_N$).}
  While operating in $\mathsf{Dec}_L$, the observer periodically runs a
  \emph{shadow collective test} on a small subset
  $A_{\mathrm{shadow}}\subseteq A_{\mathrm{mon}}$
  (with $|A_{\mathrm{shadow}}|\ll m$, requiring coherence only over the
  shadow subset):
  %
  (a)~apply a collective POVM to $A_{\mathrm{shadow}}$, obtaining
  $\widehat{x}_{\mathrm{shadow}}$;
  (b)~compare the shadow estimate's consistency with the product
  decoder's estimate over an evaluation window of length~$w$.
  %
  If the shadow collective measurement shows statistically significant
  improvement---specifically, if
  $\mathrm{MSE}_{\mathrm{shadow}} <
  (1-\theta)\,\mathrm{MSE}_{\mathrm{prod}}$
  for a contraction margin $\theta\in(0,1)$---the observer executes
  $\mathsf{Dec}_L\to\mathsf{Dec}_N$.

  \emph{Physical trigger:} equivalently, if the re-estimated
  $\widehat{T}_2$ exceeds $\tau_{\mathrm{coll}}$ for $w$ consecutive
  episodes, the return transition is warranted.
\end{enumerate}

\noindent
The monitoring fraction $f$ imposes an overhead: the effective decoding
set is reduced to $(1-f)\,m$ fragments, increasing the required total
fragment count by a factor of $1/(1-f)$ (e.g., ${\sim}12\%$ overhead
for $f=0.1$).
This overhead is the Level~2 cost
$c_{\mathrm{meta}}$ in the language of
\citep[Proposition~1]{WuDLNCompression2026preprint}.
The recovery time after coherence restoration is bounded by
$n_{\mathrm{holdout}}+w$ episodes, where $n_{\mathrm{holdout}}$ is the
minimum interval between contraction tests.


\begin{table}[t]
\caption{%
DLN instantiation in quantum Darwinism.
The belief-dependency graph~$G$ (Level~1) and revision
graph~$\mathcal{R}$ (Level~2) are the formal objects of
\citet{WuDLNCompression2026preprint};
the decoder class and error exponent are derived in this paper.
}
\label{tab:q-dln-mapping}
\centering
\small
\begin{tabular}{
  p{0.06\linewidth}
  p{0.10\linewidth}
  p{0.22\linewidth}
  p{0.22\linewidth}
  p{0.20\linewidth}
  p{0.10\linewidth}
}
\toprule
$q_O$ & DLN stage & Belief graph $G$
      & Revision graph $\mathcal{R}$
      & QD decoder class
      & Exponent \\
\midrule
$q_{\mathrm{D}}$ & Dot
  & $G=\emptyset$
  & $|\mathcal{M}|{=}1$; no edges
  & Single-fragment; memoryless
  & $0$ \\[3pt]
$q_{\mathrm{L}}$ & Linear
  & Null graph on $m$ nodes
  & Fixed point
  & Product POVM $\bigotimes_k M_k$ + classical post.
  & $-\log c$ \\[3pt]
$q_{\mathrm{N}}$ & Network
  & Bipartite DAG: $X\to\cE_k$
  & Cycle: $\mathsf{Dec}_N \rightleftharpoons \mathsf{Dec}_L$
  & Collective POVM on $\bigotimes_{k\in A}\cE_k$
  & $-\log(c^2)$ \\
\bottomrule
\end{tabular}
\end{table}


\begin{remark}[Finite-resource sample complexity]
\label{rem:finite-resource}
The fragment-count thresholds in
Corollary~\ref{cor:redundancy-stages} are leading-order
(large-$\log(1/\delta)$) statements.
For finite~$m$, the non-asymptotic bounds of
Secs.~\ref{sec:chernoff}--\ref{sec:eps-robust} (Chernoff upper bound,
calibration contraction, $\varepsilon$-robustness) apply at each DLN
stage.
This places the stage distinction in the finite-resource regime of
non-asymptotic quantum information theory
\citep{Tomamichel2016,AudenaertMosonyiVerstraete2012,
ChengDattaLiuNuradhaSalzmannWilde2025}.
\end{remark}


% =====================================================================
%  REVISED SECTION 7  (substantially shortened; core content is now in 3.3)
% =====================================================================

\section{DLN-series connection}
\label{sec:dln-brief}

The DLN instantiation developed in Sec.~\ref{sec:qstates} imports the
two formal objects of the DLN compression model---the belief-dependency
graph~$G$ and the revision graph~$\mathcal{R}$---and shows that they
specialize to the QD/SBS setting without modification:
$G$ is the conditional-independence DAG of the SBS state
(Proposition~\ref{prop:sbs-dag}), and
$\mathcal{R}$ is the directed graph of physically realizable transitions
between measurement strategies
(Definition~\ref{def:revision-qd}).

The quantitative results of this paper fill in the \emph{physics content}
of each DLN stage in the quantum setting.
The factor-of-two exponent gap
(Proposition~\ref{prop:tight-gap}) is the QD expression of the
$O(F)$-vs-$O(K)$ compression thesis of
\citep[Proposition~1(i)]{WuDLNCompression2026preprint},
and the bounded-recovery result
(Proposition~\ref{prop:recovery-qd}) is the QD expression of the
return-transition mechanism of
\citep[Proposition~1(iii)]{WuDLNCompression2026preprint}.

We emphasize three points regarding scope:
\begin{enumerate}[leftmargin=2em]
\item The quantitative claims (Chernoff sample complexity, calibration
  contraction, $\varepsilon$-robustness) are self-contained and do not
  presuppose any developmental or cognitive interpretation of DLN.
\item The DLN framework provides the \emph{structural classification}
  (which decoder topology) and the \emph{revision formalism} (when the
  observer switches); the present paper provides the \emph{quantitative
  content} (error exponents, fragment counts, robustness bounds) that
  populates each cell of that classification in a physical setting.
\item The treatment of non-stationary observation protocols---where
  $\mathcal{R}$-transitions are exercised in response to changing
  decoherence conditions---and connections to adaptive quantum
  metrology are deferred to subsequent work.
\end{enumerate}


% =====================================================================
%  REVISED ABSTRACT (last two sentences only)
%  Replace the current final sentence of the abstract with:
% =====================================================================

% CURRENT (to be removed):
%   Finally, to connect this physics formalism to the DLN observer-quality
%   thesis, we introduce a three-state stage index
%   $q_O\in\{q_{\mathrm{D}},q_{\mathrm{L}},q_{\mathrm{N}}\}$ mapping
%   Dot/Linear/Network stages to memoryless, product, and collective
%   decoding classes.

% PROPOSED REPLACEMENT:
%   Finally, we show that the conditional-independence structure of
%   SBS states is a bipartite factor DAG whose latent node is the
%   pointer variable, instantiating the belief-dependency graph of the
%   Dot--Linear--Network (DLN) compression framework without
%   modification.  This identification yields a three-stage decoder
%   classification $q_O\in\{q_{\mathrm{D}},q_{\mathrm{L}},q_{\mathrm{N}}\}$
%   with a tight, achievable factor-of-two exponent gap between product
%   and collective decoding for pure-state records, and a revision-graph
%   formalism governing adaptive transitions between measurement
%   strategies under changing coherence conditions.
